% ============================================================
\section{Discussion}
\label{sec:discussion}
% ============================================================

\subsection{Operational Sufficiency}
\label{sec:operational-sufficiency}

A deployed decision module either satisfies its requirement or it does not. The
requirement here is stated in advance and admits no partial credit: under a declared
condition, every execution names the intended action. Measured against that
requirement, an improvement in aggregate correctness is not evidence of anything in
particular. Expert role assignment produces the largest aggregate movement of the
seven conventional interventions, and converts one of the sixteen reproducibly
incorrect models. The other fifteen are exactly where they were. An aggregate figure
sums over models that changed state and models that did not, and reports a number
that belongs to neither.

The two conditions in which the aggregate points the wrong way make the objection
concrete rather than rhetorical. Encouragement yields fewer correct executions than
no intervention at all while the count of reproducibly incorrect models rises;
chain-of-thought yields more correct executions while that count stays fixed and one
model moves from inconsistent to consistent error. A practitioner reading either
figure alone would draw a conclusion the underlying states contradict.

Reproducibly correct, reproducibly incorrect and non-reproducible are therefore not
three points on a scale but three operational conditions with different consequences.
The distinction that matters is not how often a model is wrong but whether being
wrong is detectable. A non-reproducible model discloses its own unreliability to
the cheapest available check: run the input twice and the answers differ. A
reproducibly incorrect model passes that check and every repetition of it. Sixteen
of the twenty-three models in the study population fail in the second way, which
means that the standard deployment-time test certifies the majority of the fleet as
stable while they are uniformly wrong. An intervention that moves a model from the
first state to the second has made the system worse in a way no accuracy figure
will show, and four of the seven conventional interventions do exactly that.

Because the operational state depends on the input text as much as on the weights,
the unit that can carry a reliability claim is the model--condition pair rather than
the model. One model in the fleet illustrates the range. Claude Sonnet 5, on one
question, across nine wordings, is reproducibly correct under five of them,
non-reproducible under two, and reproducibly incorrect under two. The weights,
decoding configuration and retry policy are identical throughout; only the
surrounding text differs. There is no single fact about this model's reliability on
this decision to report.

A score obtained under one wording therefore establishes performance under that
wording. It does not, by itself, establish a wording-independent capability, and
treating it as one attributes to the model a property that the measurement locates
in the pair. This is not a claim that models lack capabilities, nor that all
wordings are equivalent --- the results show they are not. It is a claim about what a
single measurement licenses: the condition under which a system will be operated is
part of what is being tested, and a reliability claim that omits it is
underdetermined in the same way the untreated task itself is.
